\documentclass[12pt]{article}
\usepackage[utf8]{inputenc}
\usepackage{amsmath, amssymb, amsthm, braket}
\usepackage{graphicx}
\usepackage{geometry}
\usepackage{fancyhdr}

\geometry{letterpaper, margin=1in}
\newcommand{\Mod}[1]{\ (\mathrm{mod}\ #1)}
\setlength{\headheight}{15pt}

% Header setup
\pagestyle{fancy}
\fancyhf{}
\lhead{Math 300 Spring 2026}
\rhead{Jeffrey Lu}
\cfoot{\thepage}
\parskip=1ex

\title{Function Notes\\ \large Math 300, Spring 2026 \\ Rutgers University}
\author{Jeffrey Lu}

\begin{document}

\maketitle

\textbf{Definitions:}

Image: $f(X) = \set{f(x) \in B | x \in A}$.

Pre-Image: $f^{-1}(X) = \set{x \in X | f(x) \in B}$.

\noindent{}\textbf{Question 1:}

$f([-3,5]) = [3,28]$, so $f^{-1}([3,28]) = [-5,5]$.

\noindent{}\textbf{Question 2:}

Consider the fucntion $f: A \rightarrow B$ and a subset $X \subseteq A$.
Prove $X \subseteq f^{-1}(f(X))$.
\begin{proof}
    Suppose $f: A \rightarrow B$ and $X \subseteq A$.
    Let $x \in X$.
    By definition of image, $f(x) \subseteq f(X)$.
    By definition of preimage, $x \subseteq f^{-1}(f(X))$.
    Since $x$ is an arbitrary element of $X$, it follows that $X \subseteq f^{-1}(f(X))$.
\end{proof}

\noindent{}\textbf{Question 3:}

\begin{proof}
    Assume that $y \in f(W \cap X)$.
    Then, there exists an $x \in W \cap X$ such that $y = f(x)$.
    Since $x \in W \cap X$, it follows that $x \in W$ and $x \in X$.
    By definition of image, $f(x) \in f(W)$ and $f(x) \in f(X)$.
    Therefore $y \in f(W) \cap f(X)$, and so $f(W \cap X) \subseteq f(W) \cap f(X)$.
\end{proof}



\end{document}

